% from ega4-2-fr.tex
\oldpage[IV-2]{194}

\begin{env}[7.3.6]
\label{IV.2.s7.3.6-fr}
Nous dirons que la propriété $\boldsymbol P$ est \emph{géométrique} si elle
satisfait (en outre de $(\mathrm P_{\mathrm I})$,
$(\mathrm P_{\mathrm{II}})$ et $(\mathrm P_{\mathrm{III}})$) à la
condition :
\begin{enumerate}
\item[$(\mathrm P_{\mathrm{IV}})$] \emph{(extension de type fini du corps
de base).} --- Si $\boldsymbol P(Z,k)$ est vraie, alors, pour toute
extension $k'$ de type fini de $k$,
$\boldsymbol P(Z\otimes_k k',k')$ est vraie.
\end{enumerate}
\end{env}

\begin{lemma}[7.3.7]
\label{IV.2.s7.3.7-fr}
Soient $f:X\to Y$ un $\boldsymbol P$-morphisme de préschémas localement
noethériens, $g:Y'\to Y$ un morphisme localement de type fini. Alors, si
$\boldsymbol P$ vérifie la condition $(\mathrm P_{\mathrm{IV}})$, le
morphisme $f'=f_{(Y')}:X\times_Y Y'\to Y'$ est un
$\boldsymbol P$-morphisme.
\end{lemma}

\begin{proof}
En effet, pour tout $y'\in Y'$, si l'on pose $y=g(y')$, $k(y')$ est une
extension de type fini de $k(y)$ (I, 6.4.11) et
$f'^{-1}(y')=f^{-1}(y)\otimes_{k(y)}k(y')$ (I, 3.6.4); il suffit donc
d'appliquer $(\mathrm P_{\mathrm{IV}})$.
\end{proof}

\begin{env}[7.3.8]
\label{IV.2.s7.3.8-fr}
\emph{Exemples.} --- Les propriétés $\boldsymbol P(Z,k)$ suivantes
vérifient les conditions $(\mathrm P_{\mathrm I})$,
$(\mathrm P_{\mathrm{II}})$ et $(\mathrm P_{\mathrm{III}})$ :
\begin{enumerate}
\item[(i)] (aussi notée $(\mathrm i_n)$) $Z$ est de coprofondeur
$\leqslant n$.
\item[(ii)] $Z$ est un préschéma de Cohen--Macaulay.
\item[(iii)] (aussi notée $(\mathrm{iii}_n)$) $Z$ vérifie $(\mathrm S_n)$.
\item[(iv)] $Z$ est régulier.
\item[(v)] (aussi notée $(\mathrm v_n)$) $Z$ vérifie $(\mathrm R_n)$.
\item[(vi)] $Z$ est réduit.
\item[(vii)] $Z$ est normal.
\end{enumerate}
Pour les propriétés (ii) à (vii), cela résulte de (6.6.1) qui, en fait,
démontre la condition plus forte $(\mathrm P'_{\mathrm I})$. Pour (i), la
propriété $(\mathrm P_{\mathrm{II}})$ résulte de (6.6.2); la propriété
$(\mathrm P_{\mathrm I})$ résulte, par le même raisonnement que dans
(6.6.1, (i)), de (6.3.2) et du fait qu'un préschéma régulier est de
coprofondeur 0.

En outre, il résulte de (6.7.1) que les propriétés (i), (ii) et (iii) sont
géométriques; en vertu de (6.7.8), il en est de même des suivantes :
\begin{enumerate}
\item[(iv$'$)] $Z$ est géométriquement régulier.
\item[(v$'$)] (aussi notée $(\mathrm v'_n)$) $Z$ possède la propriété
$(\mathrm R_n)$ géométrique.
\item[(vi$'$)] $Z$ est géométriquement réduit.
\item[(vii$'$)] $Z$ est géométriquement normal.
\end{enumerate}
\end{env}

\begin{remarks}[7.3.9]
\label{IV.2.s7.3.9-fr}
\begin{enumerate}
\item[\emph{i)}] Chacune des propriétés (iv$'$), $(\mathrm v'_n)$, (vi$'$),
(vii$'$) entraîne respectivement la propriété correspondante (iv),
$(\mathrm v_n)$, (vi), (vii). La propriété (iv) implique toutes les
propriétés de (i) à (vii), et la propriété (iv$'$) implique toutes les
propriétés énumérées dans (7.3.8). Rappelons aussi que $(\mathrm i_0)$ est
équivalente à (ii); la conjonction de $(\mathrm{iii}_1)$ et de
$(\mathrm v_0)$ (resp. de $(\mathrm{iii}_1)$ et de
$(\mathrm v'_0)$) équivaut à (vi) (resp. (vi$'$)) (5.8.5); enfin la
conjonction de $(\mathrm{iii}_2)$ et $(\mathrm v_1)$ (resp. de
$(\mathrm{iii}_2)$ et $(\mathrm v'_1)$) équivaut à (vii) (resp.
(vii$'$)) (5.8.6).
\item[\emph{ii)}] Dans tous les exemples de (7.3.8), la propriété
$\boldsymbol Q(\mathscr O_z,k)$ qui sert à définir $\boldsymbol P$ est
telle que, pour toute \emph{générisation} $z'$ de $z$ dans $Z$,
$\boldsymbol Q(\mathscr O_z,k)$ \emph{entraîne}
$\boldsymbol Q(\mathscr O_{z'},k)$ : en vertu de (2.3.4), il suffit de le
vérifier pour les propriétés (i) à (vii) (et même (i) à (v) d'après la
remarque (7.3.9, (i)) : or c'est inclus dans la définition pour (iii) et
(v) (5.7.2 et 5.8.2), cela résulte de (6.3.9) pour (i), et enfin de
(0, 16.5.10 et 17.3.2) pour (ii) et (iv).
\end{enumerate}
\end{remarks}

